\section{Introduction}

Interpretability has moved from reading models to writing them. A
residual-stream direction is not only an explanation; it can serve as an
actuator. Activation steering and representation engineering make this
explicit, and recent methods cast steering as feedback or model-based control
\cite{turner2023activation,zou2023representation,nguyen2025feedback,skifstad2026local}.
This raises a basic systems question. A feedback loop has a controller and an
observer. The controller acts on a state estimate. If the estimate is a proxy
for the intended state, the loop can regulate the wrong quantity even when its
commands remain bounded.

ObserverBench makes this concern measurable. Its control task separates the
estimator from the direction along which the controller acts. We can therefore
hold a direction fixed and change only the estimate, or hold the estimate fixed
and change only the direction. The resulting factorial comparison answers a
question that the original diagonal comparison could not: whether an observed
difference comes from state-estimation error, direction geometry, or their
interaction.

We then ask the same measurement question on a real pretrained circuit. IOI is
useful because its Name Mover, Backup Name Mover, and Negative Name Mover groups
have been studied in detail \cite{wang2023ioi,mcgrath2023hydra,mcdougall2023copy,rushing2024selfrepair}.
We use those groups as borrowed structure, not as a new circuit discovery.
Whole-group ablation checks that the known conditional effect is visible.
Head-level subsets then ask whether an observer predicts held-out
interventions. This order matters: reproducing a familiar diagnostic does not
show that the same interaction controls prediction under a new intervention
distribution.

The artifact follows the same separation. Its public contract has four small
parts: a state estimator, an actuation-direction provider, a controller, and a
task runner. The task registry says which bundled fixtures accept an outside
observer. At v0, the analytic Ctl-1 task is the tested extension point. This is
enough to make the abstraction executable and falsifiable. It is not yet a
claim that every task supports arbitrary plugins.

\paragraph{What ObserverBench does.}
ObserverBench specifies an observer evaluation as a controlled measurement
experiment. The task records the target quantity, access regime, measurement
design, intervention geometry, intervention distribution, and validation
target. The validation target may be one-shot target error, collateral
movement, held-out intervention prediction, or closed-loop tracking error. The
output is an ObserverCard that says what the observer estimated, what it was
tested against, where it failed, and whether the cheap observer was sufficient.

\paragraph{What ObserverBench does not do.}
ObserverBench v0 is not a general benchmark platform, a replacement for
circuit discovery, or an adversarial oversight benchmark. It does not claim
that interaction-aware observers are always better. Several results are useful
nulls. The goal is to locate the regime in which an observer is adequate, not
to make the more complex method win.

\paragraph{Main findings.}
The results are as follows.

\begin{itemize}
  \item In Ctl-2, estimator and direction effects are separated. At every
  nonzero interaction strength, first-order integrated squared error exceeds
  lifted error in all six training seeds with either direction fixed.
  Direction geometry modulates the error. Scalar response calibration repairs
  the dynamic mismatch in the affine plant, but it leaves initial bias and can
  be ill-conditioned.

  \item In analytic and trained Ctl-1, a first-order observer can move the
  target while also moving a nuisance variable. In the analytic task, the
  result survives a base-coordinate edit that recomputes the interaction
  exactly, although the continuous edit remains a local robustness check rather
  than a realizable token intervention.

  \item In GPT-2-small IOI, random head subsets make the additive observer a
  strong baseline, while primary-stratified subsets expose cross-group
  structure. The current single-product analysis points to the Name
  Mover--Negative Name Mover pair; the capacity-matched re-analysis is treated
  as a separate experimental phase.

  \item ObserverBench exposes a versioned composition contract, deterministic
  task registry, strict result schema, and one tested external-observer adapter.
  This addresses reuse without presenting the current research code as a
  mature plugin ecosystem.
\end{itemize}
